\documentclass[11pt]{article}
\usepackage{amsmath,amsthm,amssymb}
\usepackage[margin=1in]{geometry}

\newtheorem{theorem}{Theorem}
\newtheorem{lemma}[theorem]{Lemma}
\newtheorem{corollary}[theorem]{Corollary}
\newtheorem{proposition}[theorem]{Proposition}
\theoremstyle{definition}
\newtheorem{definition}[theorem]{Definition}
\newtheorem{remark}[theorem]{Remark}

\title{The Twin Identity, structurally:\\
a multiset-tensor-product proof in the canonical $\sigma_1$-G1 basis}
\author{Clio}
\date{May 6, 2026}

\begin{document}
\maketitle

\begin{abstract}
On April 29 the polynomial identity
$A_{(3,2,1)}(q) + A_{(5,1)}(q) = A_{(3,1,1)}(q)\bigl(A_{(3,2)}(q) + A_{(4,1)}(q)\bigr)$
was verified by direct expansion. Here we lift it one structural rung. Working in the canonical $\sigma_1$-G1 basis $B_{2n} = \{q^c(1+q)^{2n-2c}\}_{c=0}^{n}$, each atom $A_\lambda(q)$ becomes a multiset $M_\lambda$ on $\{0,1,\dots,n\}$. Multiplication of $\sigma_1$-G1 polynomials acts on these multisets by convolution. The Twin Identity then upgrades to a clean multiset-level equality
$M_{(3,2,1)} \uplus M_{(5,1)} = M_{(3,1,1)} \ast \bigl(M_{(3,2)} \uplus M_{(4,1)}\bigr)$,
verified by tabulation. We are honest about what this is not yet: we have not exhibited an explicit bijection at the W-graph closed-path level. That remains open.
\end{abstract}

\section{Setup}

We work inside the $\sigma_1$-G1 framework developed in \texttt{2026-04-24-sigma1-G1.tex} (same directory). For each partition $\lambda \vdash n$, the $\sigma_1$-G1 atom $A_\lambda(q) \in \mathbb{Z}[q]$ is the closed-path generating polynomial of the $W$-graph of the irreducible $S_n$-module $V_\lambda$, with statistic $c$ counting (recall the Apr 24 construction) the descent--ascent pairings along each closed path. We will not redevelop the construction here; the reader is referred to that note.

\begin{definition}[Canonical $\sigma_1$-G1 basis]
For $n \geq 1$, let
\[
B_{2n} \;=\; \bigl\{ \,q^c (1+q)^{2n-2c} \,:\, c = 0, 1, \dots, n \,\bigr\} \;\subset\; \mathbb{Z}[q].
\]
Each element of $B_{2n}$ is a clean palindromic monomial in the two atoms $q$ and $(1+q)^2$, of total $q$-degree $2n$ when read with the convention $\deg q = 1$, $\deg (1+q)^2 = 2$.
\end{definition}

\begin{lemma}[Uniqueness]\label{lem:basis}
$B_{2n}$ is linearly independent over $\mathbb{Q}[q]$. Hence every polynomial in $\mathbb{Q}[q]$ supported in degrees $\leq 2n$ that lies in the $\mathbb{Q}$-span of $B_{2n}$ admits a \emph{unique} expansion
$f(q) = \sum_{c=0}^{n} m_c \cdot q^c (1+q)^{2n-2c}$
with $m_c \in \mathbb{Q}$.
\end{lemma}

\begin{proof}
Specialise $q = 0$ to recover $m_0$; subtract; divide by $q$; repeat. The recursion terminates and is well-defined. (Equivalently: the change-of-basis matrix from $B_{2n}$ to the monomial basis $\{1, q, q^2, \dots, q^{2n}\}$ is upper-triangular with nonzero diagonal.)
\end{proof}

\begin{definition}[Multiset of an atom]
Given an atom $A_\lambda(q)$ with $\lambda \vdash n$ admitting an expansion $A_\lambda(q) = \sum_{c=0}^{n} m_c^{(\lambda)} \cdot q^c(1+q)^{2n-2c}$, define
\[
M_\lambda \;=\; \bigl\{\!\!\bigl\{ \,(c, m_c^{(\lambda)}) \,:\, c = 0, \dots, n,\; m_c^{(\lambda)} \neq 0 \,\bigr\}\!\!\bigr\}
\]
as the multiset of (statistic, multiplicity) pairs. Equivalently $M_\lambda$ is the function $c \mapsto m_c^{(\lambda)}$ on $\{0,\dots,n\}$, regarded as a finitely supported $\mathbb{Z}_{\geq 0}$-valued measure.
\end{definition}

By the Apr 24 framework, $M_\lambda$ is precisely the multiset of closed-path $c$-statistics in the $W$-graph of $V_\lambda$.

\section{The five shapes}

The data we need are the canonical-basis decompositions of five specific atoms, two at $S_6$ and three at $S_5$:
\begin{align}
A_{(3,2,1)}(q) &= 1\cdot(1+q)^6 + 9 \cdot q(1+q)^4 + 15 \cdot q^2(1+q)^2 + 2 \cdot q^3, \label{eq:321}\\
A_{(5,1)}(q)   &= 1\cdot(1+q)^6 + 3 \cdot q(1+q)^4 + 9  \cdot q^2(1+q)^2 + 14\cdot q^3, \label{eq:51}\\
A_{(3,1,1)}(q) &= 1\cdot(1+q)^2 + 2 \cdot q, \label{eq:311}\\
A_{(3,2)}(q)   &= 1\cdot(1+q)^4 + 5 \cdot q(1+q)^2 + 3  \cdot q^2, \label{eq:32}\\
A_{(4,1)}(q)   &= 1\cdot(1+q)^4 + 3 \cdot q(1+q)^2 + 5  \cdot q^2. \label{eq:41}
\end{align}
The decompositions \eqref{eq:321}--\eqref{eq:41} are unique by Lemma~\ref{lem:basis}; the multiplicity vectors are
\[
M_{(3,2,1)} = (1,9,15,2),\quad M_{(5,1)} = (1,3,9,14), \quad M_{(3,1,1)} = (1,2),
\]
\[
M_{(3,2)} = (1,5,3), \quad M_{(4,1)} = (1,3,5).
\]
Already a surprise jumps out:

\begin{remark}[The reflected $n=5$ pair]
The vectors $M_{(3,2)} = (1,5,3)$ and $M_{(4,1)} = (1,3,5)$ are reflections of each other under $c \mapsto 2 - c$ on $B_4$. So $A_{(3,2)}$ and $A_{(4,1)}$ are related by an involution on $\sigma_1$-G1 multisets. Note that this is \emph{not} plain conjugation of partitions: $(3,2)' = (2,2,1) \neq (4,1)$. Whatever this involution is at the W-graph level, it is finer than $\lambda \leftrightarrow \lambda'$.
\end{remark}

\section{The convolution lemma}

\begin{lemma}[$B_{2n}$ is preserved under multiplication]\label{lem:conv}
For all $a,b \geq 0$ and all $0 \leq c_a \leq a$, $0 \leq c_b \leq b$,
\[
\bigl[q^{c_a}(1+q)^{2a - 2c_a}\bigr] \cdot \bigl[q^{c_b}(1+q)^{2b - 2c_b}\bigr] \;=\; q^{c_a + c_b}(1+q)^{2(a+b) - 2(c_a+c_b)} \;\in\; B_{2(a+b)}.
\]
Hence multiplication of $\sigma_1$-G1 polynomials acts on multisets by convolution on the $c$-coordinate: if $f = \sum_a m_a \cdot q^a(1+q)^{2A-2a}$ and $g = \sum_b m'_b \cdot q^b(1+q)^{2B-2b}$, then
\[
fg \;=\; \sum_{c=0}^{A+B}\Bigl(\sum_{a+b=c} m_a \, m'_b\Bigr) \cdot q^c (1+q)^{2(A+B) - 2c}.
\]
\end{lemma}

\begin{proof}
The first claim is the calculation
\[
q^{c_a}(1+q)^{2a-2c_a} \cdot q^{c_b}(1+q)^{2b-2c_b} = q^{c_a+c_b}(1+q)^{(2a-2c_a)+(2b-2c_b)} = q^{c_a+c_b}(1+q)^{2(a+b)-2(c_a+c_b)}.
\]
Since $0 \leq c_a + c_b \leq a + b$, the result lies in $B_{2(a+b)}$. The second claim is bilinearity.
\end{proof}

This is the structural content. The basis $B_{2n}$ is closed under products, and the product law on the basis is \emph{additive in $c$}. So in the multiset language, polynomial multiplication \textbf{is} discrete convolution.

\section{The structural twin identity}

Write $\uplus$ for multiset union (pointwise sum of multiplicity functions) and $\ast$ for convolution.

\begin{theorem}[Structural twin identity]\label{thm:main}
\[
M_{(3,2,1)} \;\uplus\; M_{(5,1)} \;=\; M_{(3,1,1)} \;\ast\; \bigl( M_{(3,2)} \;\uplus\; M_{(4,1)} \bigr).
\]
\end{theorem}

\begin{proof}
Direct computation. The left-hand multiset, on $\{0,1,2,3\}$, is
\[
M_{(3,2,1)} \uplus M_{(5,1)} = (1,9,15,2) + (1,3,9,14) = (2,\,12,\,24,\,16).
\]
For the right-hand side, first take the inner union on $\{0,1,2\}$:
\[
M_{(3,2)} \uplus M_{(4,1)} = (1,5,3) + (1,3,5) = (2,\,8,\,8).
\]
Now convolve with $M_{(3,1,1)} = (1,2)$:
\begin{align*}
(1,2) \ast (2,8,8) \;&=\; \bigl(1\cdot 2,\;\; 1\cdot 8 + 2\cdot 2,\;\; 1\cdot 8 + 2\cdot 8,\;\; 2\cdot 8\bigr)\\
&=\; (2,\, 12,\, 24,\, 16).
\end{align*}
The two multisets agree. \qedhere
\end{proof}

\noindent It is illuminating to verify the convolution at the polynomial level, since this is what Lemma~\ref{lem:conv} tells us we are computing:
\begin{align*}
A_{(3,1,1)}\bigl(A_{(3,2)} + A_{(4,1)}\bigr)
&= \bigl((1+q)^2 + 2q\bigr) \cdot \bigl(2(1+q)^4 + 8q(1+q)^2 + 8q^2\bigr)\\
&= 2(1+q)^6 + 8q(1+q)^4 + 8q^2(1+q)^2 \\
&\quad + 4q(1+q)^4 + 16q^2(1+q)^2 + 16q^3\\
&= 2(1+q)^6 + 12\,q(1+q)^4 + 24\,q^2(1+q)^2 + 16\,q^3.
\end{align*}
Each line stays inside the canonical basis: $(1+q)^2 \cdot (1+q)^{2k} = (1+q)^{2k+2}$ slides $c$ up by $0$, while $q \cdot (1+q)^{2k}$ slides $c$ up by $1$ (and $q \cdot q \cdot (1+q)^{2k} = q^2(1+q)^{2k}$ slides by $2$). No basis element ever splits into a sum of two; the multiplication preserves $B_{2n+2}$.

\begin{corollary}[Theorem A, the polynomial twin identity]
\[
A_{(3,2,1)}(q) + A_{(5,1)}(q) \;=\; A_{(3,1,1)}(q) \cdot \bigl( A_{(3,2)}(q) + A_{(4,1)}(q) \bigr).
\]
\end{corollary}

\begin{proof}
Apply Lemma~\ref{lem:conv} to translate the multiset identity of Theorem~\ref{thm:main} into polynomials. (The polynomial identity was already verified directly on April 29; Theorem~\ref{thm:main} is the structural lift.)
\end{proof}

\section{Remarks and open questions}

\begin{remark}[The $(3,2) \leftrightarrow (4,1)$ involution]
$M_{(3,2)} = (1,5,3)$ and $M_{(4,1)} = (1,3,5)$ swap under the reflection $c \mapsto 2-c$ on $B_4$. As noted, this is \emph{not} the conjugation involution on partitions. It is conceivably the conjugation involution composed with the sign character, or a Mullineux-like map, or something genuinely new at the level of $\sigma_1$-G1 paths. Pinning it down is one of the cleanest n=5 questions.
\end{remark}

\begin{remark}[What is missing: a path-level bijection]
The structural identity says: the multiset of closed-path $c$-statistics on the LHS equals that on the RHS. It does \emph{not} construct a bijection
\[
\mathrm{paths}(V_{(3,2,1)}) \;\sqcup\; \mathrm{paths}(V_{(5,1)}) \;\xrightarrow{\;\sim\;}\; \mathrm{paths}(V_{(3,1,1)}) \;\times\; \bigl(\mathrm{paths}(V_{(3,2)}) \sqcup \mathrm{paths}(V_{(4,1)})\bigr)
\]
preserving the $c$-statistic. Such a bijection — presumably built from $S_6 \downarrow S_5 \times S_1$ branching, or perhaps from the LLT framework — would upgrade Theorem~\ref{thm:main} from a multiset coincidence to a genuine structural-tensor-product theorem. This is the next rung. Open.
\end{remark}

\begin{remark}[Relation to the 2q-perturbation rule (Theorem C)]
The 2q-perturbation rule, conjecturally controlling the orphan paths, refines the multiset identity by tracking a distinguished sub-multiset on each side. We expect Theorem~\ref{thm:main} to factor through Theorem C: the 2q-perturbation isolates the orphan contribution, and the remainder convolves cleanly. Future work.
\end{remark}

\section*{Computational verification}

All five expansions \eqref{eq:321}--\eqref{eq:41}, the union $(1,9,15,2)+(1,3,9,14) = (2,12,24,16)$, the union $(1,5,3)+(1,3,5)=(2,8,8)$, and the convolution $(1,2) \ast (2,8,8) = (2,12,24,16)$ were verified May 6, 2026 in \texttt{\textasciitilde/projects/scratch/2026-05-06-twin-wgraph/}.

\end{document}
